\documentclass[11pt]{article}
\usepackage[margin=1in]{geometry}
\usepackage{amsmath,amssymb,amsthm}
\usepackage[colorlinks=true,linkcolor=blue,citecolor=blue,urlcolor=blue]{hyperref}

\newtheorem{theorem}{Theorem}
\newtheorem{corollary}{Corollary}

\newcommand{\R}{\mathbb R}
\newcommand{\T}{\mathbb T}
\newcommand{\Z}{\mathbb Z}
\newcommand{\F}{\mathcal F}

\begin{document}

\begin{center}
{\Large A planar reduction for real frame connectedness}\\[4pt]
{\large arXiv:2108.02275 residual real-frame problem}
\end{center}

\paragraph{Status.}
Partial result, likely valid. This is an exact \(d=2\) reduction/characterization; it is not a full solution of the real-frame connectedness problem in all dimensions.

\paragraph{Source problem.}
Needham--Shonkwiler, \emph{Admissibility and Frame Homotopy for Quaternionic Frames}, arXiv:2108.02275, proves the quaternionic admissibility and homotopy questions. In the introduction the authors note that, for real frames with nonconstant spectrum or vector norms, the analogue remains open. In the discussion they state that some real spaces \(\F^{\R^d,N}_{\lambda}(r)\) are connected and others are not, and that it is challenging to characterize which data \((\lambda,r)\) lead to which outcome.

\paragraph{Definitions.}
For \(r=(r_1,\dots,r_N)\in(0,\infty)^N\) and \(\lambda=(\lambda_1,\lambda_2)\), \(\lambda_1\geq\lambda_2>0\), write
\[
\F^{\R^2,N}_{\lambda}(r)
=
\left\{
F=[f_1|\cdots|f_N]\in\R^{2\times N}:
\|f_i\|^2=r_i,\ \operatorname{spec}(FF^T)=\{\lambda_1,\lambda_2\}
\right\}.
\]
Assume the necessary trace identity
\[
\sum_{i=1}^N r_i=\lambda_1+\lambda_2,
\]
and set \(\delta=\lambda_1-\lambda_2\).

\paragraph{Proof Intuition.}
In the plane, a vector of fixed length is determined by an angle \(\theta_i\). The frame operator only sees the rank-one matrices \(u_{\theta_i}u_{\theta_i}^T\), and their traceless parts rotate at double angle \(2\theta_i\). Thus fixing the two eigenvalues of the frame operator is the same as fixing the modulus of the complex phasor sum \(\sum_i r_i e^{2i\theta_i}\). The real sign ambiguity of each vector is exactly the coordinatewise double cover \(\theta_i\mapsto 2\theta_i\), so connectedness is controlled by ordinary labelled planar polygon spaces together with a finite mod-2 winding calculation.

\begin{theorem}[Planar frame-polygon double-cover reduction]
Let \(r_i>0\), \(\lambda_1\geq\lambda_2>0\), \(\sum_i r_i=\lambda_1+\lambda_2\), and \(\delta=\lambda_1-\lambda_2\). Define
\[
\Theta(r,\delta)
=
\left\{\theta=(\theta_1,\dots,\theta_N)\in\T^N:
\left|\sum_{i=1}^N r_i e^{2i\theta_i}\right|=\delta
\right\}.
\]
Then
\[
\theta\longmapsto
\left[\sqrt{r_1}(\cos\theta_1,\sin\theta_1)\ |\ \cdots\ |\ \sqrt{r_N}(\cos\theta_N,\sin\theta_N)\right]
\]
is a homeomorphism from \(\Theta(r,\delta)\) onto \(\F^{\R^2,N}_{\lambda}(r)\).

Let
\[
L(r,\delta)
=
\left\{\phi=(\phi_1,\dots,\phi_N)\in\T^N:
\left|\sum_{i=1}^N r_i e^{i\phi_i}\right|=\delta
\right\},
\]
and let \(D:\T^N\to\T^N\) be \(D(\theta_1,\dots,\theta_N)=(2\theta_1,\dots,2\theta_N)\). Under the above homeomorphism,
\[
\F^{\R^2,N}_{\lambda}(r)\cong D^{-1}(L(r,\delta)).
\]
For each path component \(C\) of \(L(r,\delta)\), the number of path components of \(\F^{\R^2,N}_{\lambda}(r)\) lying over \(C\) is
\[
\left[(\Z/2)^N:H_C\right],
\]
where \(H_C\) is the image of the composite
\[
H_1(C;\Z/2)\longrightarrow H_1(\T^N;\Z/2)\cong(\Z/2)^N.
\]
\end{theorem}

\begin{proof}
Every frame in \(\F^{\R^2,N}_{\lambda}(r)\) has nonzero columns, so it can be written uniquely by angles \(\theta_i\in\T\) as
\[
f_i=\sqrt{r_i}(\cos\theta_i,\sin\theta_i).
\]
For such a frame,
\[
FF^T
=
\sum_{i=1}^N r_i
\begin{pmatrix}
\cos^2\theta_i&\cos\theta_i\sin\theta_i\\
\cos\theta_i\sin\theta_i&\sin^2\theta_i
\end{pmatrix}.
\]
Using the double-angle identities, this becomes
\[
FF^T
=
\frac{R}{2}I_2+
\frac12
\begin{pmatrix}
A&B\\
B&-A
\end{pmatrix},
\qquad
R=\sum_i r_i,
\]
where
\[
A=\sum_i r_i\cos(2\theta_i),
\qquad
B=\sum_i r_i\sin(2\theta_i).
\]
The eigenvalues of this matrix are
\[
\frac{R}{2}\pm\frac12\sqrt{A^2+B^2}
=
\frac{R}{2}\pm\frac12\left|\sum_i r_i e^{2i\theta_i}\right|.
\]
Since \(R=\lambda_1+\lambda_2\), the spectrum is \(\{\lambda_1,\lambda_2\}\) exactly when
\[
\left|\sum_i r_i e^{2i\theta_i}\right|=\lambda_1-\lambda_2=\delta.
\]
This proves the first homeomorphism. The second identification is just the same condition after applying the coordinatewise double cover \(D\).

Finally, \(D:\T^N\to\T^N\) is the regular covering with deck group \((\Z/2)^N\). Restrict it to any path component \(C\) of \(L(r,\delta)\). By standard covering theory, the number of path components in \(D^{-1}(C)\) is the index of the monodromy subgroup of \((\Z/2)^N\). For the torus double cover, this monodromy subgroup is exactly the mod-2 winding image of \(H_1(C;\Z/2)\) inside \(H_1(\T^N;\Z/2)\cong(\Z/2)^N\). This gives the component formula.
\end{proof}

\begin{corollary}[Polygon-space form]
If \(\delta>0\), then \(L(r,\delta)\) is homeomorphic to \(S^1\times P(r,\delta)\), where
\[
P(r,\delta)=
\left\{\beta\in\T^N:\sum_{i=1}^N r_i e^{i\beta_i}=\delta\right\}.
\]
Equivalently, \(P(r,\delta)\) is the labelled planar polygon space with side lengths \(r_1,\dots,r_N,\delta\) and the \(\delta\)-side fixed on the real axis. If \(\delta=0\), \(L(r,0)\) is the usual closed \(N\)-gon phasor level.
\end{corollary}

\begin{proof}
For \(\delta>0\), the phasor sum never vanishes. For \(\phi\in L(r,\delta)\), let \(\alpha=\arg\sum_i r_i e^{i\phi_i}\) and set \(\beta_i=\phi_i-\alpha\). Then \(\sum_i r_i e^{i\beta_i}=\delta\). This gives a homeomorphism \(L(r,\delta)\to S^1\times P(r,\delta)\), with inverse \((\alpha,\beta)\mapsto(\alpha+\beta_i)_i\). The \(\delta=0\) statement is the definition of a closed phasor polygon.
\end{proof}

\paragraph{Limitations.}
The theorem reduces the entire planar real-frame problem to labelled planar polygon topology plus a finite monodromy computation. It does not solve the higher-dimensional real case, where the source obstruction is the multiplicity-one isoparametric setting rather than the double-angle phenomenon. It also does not assert that all admissible planar real frame spaces are connected; the component formula can detect disconnected examples.

\paragraph{Novelty check.}
The cheap run indexes were searched for arXiv:2108.02275 and for frame-homotopy, fixed-spectrum, fixed-norm, and polygon-space keywords. Web searches on 2026-07-03 found the source paper, the Cahill--Mixon--Strawn finite unit norm tight frame connectivity paper, and the Needham--Shonkwiler complex symplectic-frame paper. I did not find a later full characterization of the source's nonconstant real-frame problem or this exact arbitrary-\((\lambda,r)\) planar double-cover component formula. The reduction is close in spirit to the known correspondence between real planar FUNTFs and closed polygonal chains, but it includes arbitrary planar spectra and arbitrary vector norms.

\paragraph{Human-review recommendation.}
Review the convention that frame vectors are ordered and signs are not quotiented. Under that convention, the double cover and the mod-2 component formula are the essential content. If a reviewer instead works modulo independent sign changes of frame vectors, the finite covering layer disappears and the answer is just the labelled polygon level \(L(r,\delta)\).

\paragraph{References.}
\begin{itemize}
\item T. Needham and C. Shonkwiler, \emph{Admissibility and Frame Homotopy for Quaternionic Frames}, arXiv:2108.02275.
\item J. Cahill, D. G. Mixon, and N. Strawn, \emph{Connectivity and Irreducibility of Algebraic Varieties of Finite Unit Norm Tight Frames}, arXiv:1311.4748.
\item T. Needham and C. Shonkwiler, \emph{Symplectic Geometry and Connectivity of Spaces of Frames}, arXiv:1804.05899.
\end{itemize}

\end{document}
